% topic: algorithms
\question ร้านอาหารในโรงเรียนชายล้วนแห่งหนึ่งมีระบบการเข้าคิวซื้ออาหารที่แปลกประหลาด นักเรียนแต่ละคนจะมีค่าความเก๋า k และมีกฎที่จะทำงานทุก 1 นาทีดังนี้
\begin{enumerate}
    \item ถ้านาทีนี้มีคนมาเข้าแถว คนนั้นจะไล่ทุกคนในแถวที่มีค่าความเก๋าน้อยกว่าแต่อยู่ข้างหน้าเขาออกจากแถว จนกว่าจะเจอคนที่มีค่าความเก๋ามากกว่าหรือเท่ากันอยู่ข้างหน้าหรืออยู่หัวแถว
    \item หลังจากนั้น คนหน้าสุดจะได้รับอาหารและออกจากแถว
\end{enumerate}
โดยในวันนี้มีเหตุการณ์เกิดขึ้นดังนี้ ก่อนที่ร้านจะเปิดตอน 9 โมงตรงมีคนเข้าแถวดังนี้
\begin{enumerate}
    \item นาย A, k = 100 (หัวแถว)
    \item นาย B, k = 80
    \item นาย C, k = 50
    \item นาย D, k = 40 (ปลายแถว)
\end{enumerate}
และเมื่อร้านเปิดตอน 9 โมงตรง มีคนเข้าแถวตามเวลาดังนี้
\begin{enumerate}
    \item นาย E, k = 50, 09:01
    \item นาย F, k = 40, 09:02
    \item นาย G, k = 90, 09:03
    \item นาย D, k = 30, 09:04
\end{enumerate}
โดยคนแรกจะได้อาหารตอน 09:00 คนที่สองจะได้อาหารตอน 09:01 ใครได้อาหารบ้างตามลำดับจากก่อนไปหลัง \\
\choiceshorizontal{A B C G D}{A B C G}{A B E G D}{A B C E G D}
% answer: 1